% c05 —— §6 末 + §7 前半（p011–p013 / 原书 p.137–139）
% 承 c04：原书 p.136 末句「局所系 $H^{amb}_{A,\mathbb{Z}}$ のファイバーは…数値」跨页，整句归本块，
% 从「局部系统 $H^{amb}_{A,\mathbb{Z}}$ 的纤维，与…同构」起写。
% 本块止于「……这给出各纤维上 3 阶的拐点。」；原书 p.139 末句「Hesse 群 $G_{216}\cong\cdots$ は
% $S(3)\to X(3)$ に同変に作用する」跨页，整句归 c06。
% 图：图 7.1（src/figures/fig7_1.png）。

% ===== 原书 p.137 =====

局部系统 $H^{amb}_{A,\mathbb{Z}}$ 的纤维，与由从 $\mathcal{N}(\check{X})$ 拉回生成的、数值 Grothendieck
群 $\mathcal{N}(\check{Y})$ 的子群 $\mathcal{N}^{amb}(\check{Y})$ 同构。由 Riemann–Roch 定理，
$\mathcal{N}(\check{Y})$ 同构于 Grothendieck 群 $K(\check{Y})$ 在映射
\begin{equation}
  v\colon K(\check{Y})\to H^{*}(\check{Y},\mathbb{Q}),\quad
  \mathcal{E}\mapsto\ch(\mathcal{E})\cup\hat{\Gamma}_{\check{Y}} \tag{6.2}
\end{equation}
下的像。这里，$\hat{\Gamma}$ 类为 Todd 类的平方根 [Iri09]。

% [Tutor Note: (6.3) 左端的符号是 $\varsigma$（希腊字母 final sigma），不是 $\sigma$：原书图版中该字形开口、无闭合圈，与同书 (4.17) 的 $\sigma$（闭合碗形）明显不同；文字层把两者都提成 \sigma，不可据以判读。]
\begin{theorem}[Iritani {[Iri11, Theorem 6.9]}]
  存在纯极化整 Hodge 结构的变形的同构
  \begin{equation}
    \mathrm{Mir}_{\mathcal{Y}}\colon
    \varsigma^{*}(H^{amb}_{A,\mathbb{Z}},\nabla^{A},\mathscr{F}^{\bullet}_{A},Q_{A})
    \xrightarrow{\sim}(H^{vc}_{B,\mathbb{Z}},\nabla^{B},\mathscr{F}^{\bullet}_{B},Q_{B}) \tag{6.3}
  \end{equation}
\end{theorem}

定理 6.1 的证明有一个重要步骤，即 [Iri11, Theorem 5.7]；在其证明过程中证明了：将极大退化点附近
$\pi_{1}(\mathbb{L}^{\vee}_{\mathbb{C}^{\times}})\cong\mathbb{L}^{\vee}$ 的元素 $\ell$ 沿其方向的
$H^{vc}_{B,\mathbb{Z}}$ 单值性经同构 (6.3) 搬运后，它由满足
$c_{1}(\mathcal{L}_{i})=\ell\in\mathbb{L}^{\vee}\cong\Pic\check{X}$ 的 $\check{X}$ 上的线丛
$\mathcal{L}_{i}$ 的对偶丛的限制作张量积所诱导的等距映射
\begin{equation}
  (-)\otimes\iota^{*}(\mathcal{L}_{i}^{\vee})\colon
  \mathcal{N}(\check{Y})\to\mathcal{N}(\check{Y}) \tag{6.4}
\end{equation}
给出。注意，该作用可提升为相干层的导出范畴 $D^{b}\operatorname{coh}\check{Y}$ 的自等价。周期映射的单值性与
镜像侧相干层的导出范畴的自等价之间的关系可追溯到 [Kon98, Hor05]。

当 $\check{Y}$ 为 K3 曲面时，数值 Grothendieck 群 $\mathcal{N}(\check{Y})$ 分解为 Picard 群
$\Pic(\check{Y})$ 与由结构层和点层的类生成的自由阿贝尔群的直和
$\mathbb{Z}[\mathcal{O}_{\check{Y}}]\oplus\Pic(\check{Y})\oplus\mathbb{Z}[\mathcal{O}_{p}]$。
这里，Picard 群到数值 Grothendieck 群的嵌入由 $[\mathcal{O}_{\check{Y}}(D)]\mapsto[\mathcal{O}_{D}]$
给出，其上的配对为
\begin{gather}
  (\mathcal{O}_{\check{Y}},\mathcal{O}_{\check{Y}})=-2,\quad
  (\mathcal{O}_{\check{Y}},\mathcal{O}_{D})=-\chi(\mathcal{O}_{D}),\quad
  (\mathcal{O}_{\check{Y}},\mathcal{O}_{p})=-1, \tag{6.5}\\
  (\mathcal{O}_{D},\mathcal{O}_{E})=D\cdot E,\quad
  (\mathcal{O}_{D},\mathcal{O}_{p})=-\chi(\mathcal{O}_{p},\mathcal{O}_{p})=0 \tag{6.6}
\end{gather}
这里 $\mathcal{O}_{\check{Y}}$ 是结构层，$\mathcal{O}_{D}$ 是除子 $D$ 的结构层，而
$\mathcal{O}_{p}$ 是点的结构层。自等价
$(-)\otimes\mathcal{O}_{\check{Y}}(-D)\colon K(\check{Y})\to K(\check{Y})$ 的作用由
\begin{align}
  &[\mathcal{O}_{\check{Y}}]\mapsto[\mathcal{O}(-D)]=[\mathcal{O}_{\check{Y}}]-[\mathcal{O}_{D}], \tag{6.7}\\
  &[\mathcal{O}_{E}]\mapsto[\mathcal{O}_{E}(-D)]=[\mathcal{O}_{E}]-(D\cdot E)[\mathcal{O}_{p}], \tag{6.8}\\
  &[\mathcal{O}_{p}]\mapsto[\mathcal{O}_{p}] \tag{6.9}
\end{align}
给出。

% ===== 原书 p.138 =====

\section{模曲线 $X_{0}(3)$ 上的椭圆曲线族}

作为二次栈的例子，考虑 $A=\{v_{0},v_{1},v_{2},v_{3}\}$ 且
\begin{equation}
  (v_{0}\ v_{1}\ v_{2}\ v_{3})=\begin{pmatrix}0&1&0&-1\\0&0&1&-1\end{pmatrix} \tag{7.1}
\end{equation}
成立的情形。将 $\mathbb{R}^{A}$ 的元素用满足 $f_{i}(e_{j})=\delta_{ij}$ 的函数族
$\{f_{i}\}_{i=0}^{3}$ 与 $(\lambda_{i})_{i=0}^{3}\in\mathbb{R}^{4}$ 写成
$f=\sum_{i=0}^{3}\lambda_{i}f_{i}$。由 $f$ 诱导的 $\Delta$ 的多面体分割为
\begin{equation}
  \underline{\Delta}=
  \begin{cases}
    \underline{\Delta}_{-}=\{\Delta\} & 3\lambda_{0}\ge\lambda_{1}+\lambda_{2}+\lambda_{3},\\
    \underline{\Delta}_{+}=\{\Delta_{1},\Delta_{2},\Delta_{3}\} & 3\lambda_{0}<\lambda_{1}+\lambda_{2}+\lambda_{3},
  \end{cases} \tag{7.2}
\end{equation}
二次扇由两个锥 $C(\{\Delta\})=\{3\lambda_{0}\ge\lambda_{1}+\lambda_{2}+\lambda_{3}\}$ 与
$C(\{\Delta_{1},\Delta_{2},\Delta_{3}\})=\{3\lambda_{0}\le\lambda_{1}+\lambda_{2}+\lambda_{3}\}$ 组成。
多面体分割 $\underline{\Delta}_{-}$ 具有 1 个面 $\{\Delta\}$、3 条边
$\{\Delta_{10},\Delta_{20},\Delta_{30}\}$ 以及 3 个顶点 $\{v_{1},v_{2},v_{3}\}$。

\begin{figure}[htbp]
  \centering
  \includegraphics[width=0.45\textwidth]{fig7_1.png}
  \caption{$\Delta=\Conv A$ 的多边形分割（原书 p.138）}
\end{figure}

取 $\mathbb{R}^{A}/\mathbb{R}$ 的坐标为 $\overline{\lambda}_{i}=\lambda_{i}-\lambda_{0}$，
$i=1,2,3$，则 Lafforgue 扇的锥由
\begin{gather*}
  C(\underline{\Delta}_{-},\Delta)=\{\overline{\lambda}_{1}=\overline{\lambda}_{2}=\overline{\lambda}_{3}\le 0\},\\
  C(\underline{\Delta}_{-},\Delta_{10})=\{\overline{\lambda}_{2}=\overline{\lambda}_{3}\le\min\{\overline{\lambda}_{1},0\},\overline{\lambda}_{1}+\overline{\lambda}_{2}+\overline{\lambda}_{3}\le 0\},\\
  C(\underline{\Delta}_{-},v_{1})=\{\overline{\lambda}_{1}\le\min\{\overline{\lambda}_{2},\overline{\lambda}_{3},0\},\overline{\lambda}_{1}+\overline{\lambda}_{2}+\overline{\lambda}_{3}\le 0\},
\end{gather*}
等给出。

% ===== 原书 p.139 =====

多面体分割 $\underline{\Delta}_{+}$ 由 3 个面 $\{\Delta_{1},\Delta_{2},\Delta_{3}\}$、6 条边
$\{\Delta_{12},\ldots,\Delta_{30}\}$ 以及 4 个顶点 $\{v_{0},v_{1},v_{2},v_{3}\}$ 组成，
其 Lafforgue 扇的锥由
% [Tutor Note: 本段讨论的是 $\underline{\Delta}_{+}$ 的 Lafforgue 扇锥，但原书这五式的第一个参数
%  均印作 $\underline{\Delta}_{-}$；按原图照录，疑为原书笔误（对照 (7.2)，$\underline{\Delta}_{+}$ 的加号清晰可见）。]
\begin{gather*}
  C(\underline{\Delta}_{-},\Delta_{1})=\{\overline{\lambda}_{2}=\overline{\lambda}_{3}=0,\overline{\lambda}_{1}\ge 0\},\\
  C(\underline{\Delta}_{-},\Delta_{10})=\{\overline{\lambda}_{2}=\overline{\lambda}_{3}\le 0,\overline{\lambda}_{1}+\overline{\lambda}_{2}+\overline{\lambda}_{3}\ge 0\},\\
  C(\underline{\Delta}_{-},\Delta_{23})=\{\overline{\lambda}_{1}=0,\overline{\lambda}_{2}\ge 0,\overline{\lambda}_{3}\le 0\},\\
  C(\underline{\Delta}_{-},v_{0})=\{\min\{\overline{\lambda}_{1},\overline{\lambda}_{2},\overline{\lambda}_{3}\}\le 0\},\\
  C(\underline{\Delta}_{-},v_{1})=\{\overline{\lambda}_{1}\le\min\{\overline{\lambda}_{2},\overline{\lambda}_{3},0\},\overline{\lambda}_{1}+\overline{\lambda}_{2}+\overline{\lambda}_{3}\ge 0\},
\end{gather*}
等给出。整体来看，Lafforgue 扇 $\tilde{\mathcal{F}}(A)$ 有 7 个三维锥。其中 3 个对应于
$\mathbb{P}^{2}/(\mathbb{Z}/3\mathbb{Z})$ 的环面不动点，4 个对应于 3 个 $\mathbb{P}^{2}$ 的并的环面不动点。
Laurent 多项式 (5.6) 为
\begin{equation}
  W=a_{0}+a_{1}x+a_{2}y+\frac{a_{3}}{xy}, \tag{7.3}
\end{equation}
其判别式为
\[
  \Delta=a_{0}^{3}+27a_{1}a_{2}a_{3}.
\]
二次多面体为
\[
  \Sigma(A)=\Conv\{(3,0,0,0),(0,1,1,1)\},
\]
其法扇确实与二次扇 $\mathcal{F}(\Sigma)$ 一致。二次栈的稠密环面的坐标由
\begin{equation}
  v=\frac{a_{1}a_{2}a_{3}}{a_{0}^{3}} \tag{7.4}
\end{equation}
给出。

三次曲线的 Hesse 铅笔由
\begin{equation}
  X^{3}+Y^{3}+Z^{3}+\mu XYZ=0 \tag{7.5}
\end{equation}
给出。该铅笔的基点由 9 个拐点组成，在这里爆破便得到盐田的椭圆模曲面 $S(3)\to X(3)$ [Shi72]。
这里，$X(3)$ 是在 $\mathbb{H}/\Gamma(3)$ 上添加 4 个尖点所得的模曲线，每个尖点之上都带有
$I_{3}$ 型奇异纤维。9 条例外曲线给出 9 个截面，后者在每个纤维上给出 3 阶拐点。

% CHUNK c05 DONE pages p011–p013
